\documentclass[11pt,a4paper]{article}
\usepackage{fontspec}
\usepackage[french]{babel}
\usepackage[margin=2.4cm]{geometry}
\usepackage{amsmath,amssymb}
\usepackage{graphicx}
\usepackage{booktabs}
\usepackage[font=small,labelfont=bf]{caption}

\graphicspath{{../rapport/figures/}}
\newcommand{\ft}[1]{\widehat{#1}}
\setlength{\parskip}{0.4em}

\title{L'expérience sur l'IRM 3D\\[0.2em]\large Ce qui a été fait, pour pouvoir l'expliquer à l'oral}
\author{Célian Baffalie}
\date{}

\begin{document}
\maketitle

Dans le rapport, cette partie est réduite à une figure et un paragraphe. Ce document détaille tout ce qui a
été fait, dans l'ordre du code (section 10 de \texttt{docs/rapport/generer\_figures.py}), et les questions
qu'on peut me poser.

\section{Le but}

Montrer que la méthode de l'article marche aussi en 3D, sur une vraie image médicale, et comparer deux
façons de traiter un volume :
\begin{itemize}
  \item \textbf{en 3D} : on résout une seule équation de Poisson sur tout le volume, avec une FFT 3D
  (fonction \texttt{contrast\_enhance\_3d}) ;
  \item \textbf{coupe par coupe} : on applique la version 2D (\texttt{contrast\_enhance}) à chaque coupe
  axiale séparément, comme si les coupes étaient des images indépendantes.
\end{itemize}
Le fantôme (volume synthétique) a servi à vérifier le solveur 3D et à mesurer les effets ; l'IRM sert à
voir si les conclusions tiennent sur une vraie image.

\section{Les données}

\begin{itemize}
  \item \textbf{Source} : jeu de données public OpenNeuro \texttt{ds000001}, sujet \texttt{sub-01},
  image \texttt{T1w} (Schonberg et al., 2012, licence CC0). C'est une IRM anatomique d'un sujet sain.
  \item \textbf{Format} : NIfTI (\texttt{.nii.gz}), le format standard des IRM en recherche. Je le lis
  avec la bibliothèque \texttt{nibabel} : \texttt{nib.load(...)} donne le volume et un en-tête.
  \item \textbf{Taille} : $160\times192\times192$ voxels. L'en-tête donne la taille d'un voxel :
  $1\times1{,}33\times1{,}33$\,mm (fonction \texttt{get\_zooms}). Les voxels ne sont donc pas cubiques.
  \item \textbf{Axes} : le premier axe va de gauche à droite (R), le deuxième de l'arrière vers l'avant
  (A), le troisième du bas vers le haut (S).
  \item \textbf{Normalisation} : les valeurs d'une IRM n'ont pas d'unité fixe. Je les ramène à peu près
  dans $[0,255]$ en divisant par le percentile 99,5 et en multipliant par 255, pour que le seuil $T$ ait
  le même sens que sur les photos. Le percentile plutôt que le maximum évite qu'un seul voxel très clair
  fixe l'échelle.
\end{itemize}

\paragraph{Ce qu'on voit sur une IRM T1.} La pondération T1 donne : substance blanche claire, substance
grise intermédiaire, liquide cérébrospinal (ventricules, sillons) sombre, air noir, graisse (sous la peau,
moelle osseuse) très claire, os cortical sombre. Le bloc noir devant le visage vient de l'anonymisation :
le visage a été effacé pour qu'on ne puisse pas reconnaître le sujet (\emph{defacing}).

\begin{figure}[h]
\centering
\includegraphics[width=\linewidth]{irm_3d.png}
\caption{L'IRM d'origine (figure retirée du rapport). (a, b, c) Les trois coupes orthogonales.
(d) Les trois coupes placées dans l'espace.}
\end{figure}

\section{Afficher un volume : les trois coupes}

Un volume ne s'affiche pas directement ; on montre des coupes planes :
\begin{itemize}
  \item \textbf{sagittale} : on fixe le premier indice, \texttt{v[i]} (vue de profil) ;
  \item \textbf{coronale} : on fixe le deuxième, \texttt{v[:, j]} (vue de face) ;
  \item \textbf{axiale} : on fixe le troisième, \texttt{v[:, :, k]} (vue de dessus).
\end{itemize}
\texttt{np.rot90} tourne chaque coupe pour avoir le haut de la tête en haut. Comme les voxels ne sont pas
cubiques, j'affiche chaque coupe avec un rapport d'aspect égal au rapport des tailles de voxels
(\texttt{aspect} dans \texttt{imshow}), sinon la tête serait déformée. La vue 3D (d) place les trois
coupes dans l'espace avec \texttt{plot\_surface} de Matplotlib : chaque coupe est une surface plane
coloriée avec ses niveaux de gris.

\section{Pourquoi réduire le volume d'un facteur 2}

Le solveur travaille sur le volume prolongé par miroir, deux fois plus grand dans chaque direction :
$320\times384\times384\approx 47$ millions de voxels. Un tableau de nombres complexes (16 octets par
valeur) fait alors 755\,Mo. Or \texttt{contrast\_enhance\_3d} garde en mémoire plusieurs tableaux de cette
taille en même temps : $\ft f$, les trois composantes du gradient en Fourier, les trois composantes en
espace, leurs trois FFT, puis $\ft u$. Cela fait de l'ordre de 6\,Go, trop pour mon ordinateur.

J'ai donc réduit le volume en faisant la moyenne de chaque bloc de $2\times2\times2$ voxels :
\texttt{v.reshape(80, 2, 96, 2, 96, 2).mean(axis=(1, 3, 5))}. Le volume passe à $80\times96\times96$,
le volume prolongé à 5,9 millions de voxels (94\,Mo par tableau), et le calcul prend 5 à 6\,s.

\section{L'éq.~15 en 3D}

C'est exactement la méthode 2D, avec un axe de plus :
\begin{enumerate}
  \item prolongement miroir sur les trois axes, FFT 3D ;
  \item gradient en Fourier : $\ft{V_1} = i w_m \ft f$, $\ft{V_2} = i w_n \ft f$, $\ft{V_3} = i w_p \ft f$ ;
  \item retour en espace, multiplication par $\alpha$ des trois composantes dans la zone
  $X=\{f<T\}$ ;
  \item retour en Fourier et résolution :
  \[
    \ft u_{mnp} = \frac{-i\,(w_m \ft{V_1} + w_n \ft{V_2} + w_p \ft{V_3})}{w_m^2 + w_n^2 + w_p^2},
    \qquad \ft u_{000} = \ft f_{000} ;
  \]
  \item FFT inverse et on garde le huitième $H\times W\times D$ de départ.
\end{enumerate}
J'ai pris $T=60$ et $\alpha=2{,}5$. Avec $T=60$, la zone sombre contient l'air, le liquide cérébrospinal
et l'os cortical : 75\,\% du volume réduit, surtout de l'air autour de la tête.

\section{3D contre coupe par coupe}

\begin{figure}[h]
\centering
\includegraphics[width=0.62\linewidth]{irm_contraste_3d.png}
\caption{La figure gardée dans le rapport. Ligne 1 : original réduit. Ligne 2 : éq.~15 en 3D. Ligne 3 :
éq.~15 coupe axiale par coupe axiale. Les deux résultats sont étirés en min-max.}
\end{figure}

\paragraph{Ce qu'on voit.} Coupe par coupe, des stries horizontales apparaissent sur les coupes sagittale
et coronale. Elles n'existent pas en 3D, ni sur les coupes axiales.

\paragraph{Pourquoi.} Coupe par coupe, chaque coupe axiale est traitée seule :
\begin{itemize}
  \item rien ne relie une coupe à la suivante : le solveur 2D ne connaît pas la direction $z$ ;
  \item chaque coupe reçoit sa propre constante : $\ft u_{00}=\ft f_{00}$ impose que \emph{chaque coupe}
  garde sa moyenne, et le décalage dû au contraste (le saut supplémentaire $E$ du rapport) dépend de la
  proportion de zone sombre dans la coupe, qui varie d'une coupe à l'autre.
\end{itemize}
Deux coupes voisines reçoivent donc des décalages différents. Sur une coupe axiale on ne le voit pas (tout
le plan est décalé pareil), mais sur une coupe sagittale ou coronale, qui traverse toutes les coupes
axiales, chaque ligne horizontale correspond à une coupe axiale, d'où les stries. En 3D, la direction $z$
est traitée comme les deux autres et il n'y a qu'une constante pour tout le volume.

\paragraph{La mesure.} Pour chiffrer ce que l'on voit, je compare la dérivée selon $z$ du résultat avec
celle que demande l'éq.~15 :
\[
  V_z = \alpha^{[f<T]}\,\big(f[\cdot,\cdot,k+1]-f[\cdot,\cdot,k]\big),
  \qquad
  \text{écart} = \operatorname{moyenne}\,\big|\,(x[\cdot,\cdot,k+1]-x[\cdot,\cdot,k]) - V_z\,\big|,
\]
pour $x$ = l'original, le résultat 3D et le résultat coupe par coupe :

\begin{center}
\begin{tabular}{@{}lc@{}}
\toprule
 & Écart à l'éq.~15 selon $z$ \\
\midrule
Original (rien n'est fait) & 5,25 \\
Éq.~15 en 3D & 5,38 \\
Éq.~15 coupe par coupe & 9,60 \\
\bottomrule
\end{tabular}
\end{center}

\paragraph{Comment le lire.} Coupe par coupe, l'écart est presque deux fois plus grand qu'en 3D : les
coupes voisines ne sont plus cohérentes. En 3D, l'écart reste du même ordre que pour l'original : le
résultat 3D ne réalise pas exactement le champ demandé, parce que ce champ ($\alpha\nabla f$ dans $X$,
$\nabla f$ ailleurs) n'est pas le gradient d'une image. La méthode trouve l'image dont le gradient est le
plus proche au sens des moindres carrés, en répartissant l'erreur dans les trois directions. Cette mesure
utilise des différences finies alors que le solveur utilise le gradient spectral : c'est un indicateur
pour comparer les deux méthodes entre elles, pas une erreur exacte.

\section{L'étirement}

Le résultat 3D va de $-153$ à $680$, celui coupe par coupe de $-149$ à $692$. L'étirement min-max
multiplie donc les contrastes par $255/833\approx0{,}31$ (et 0,30) : l'image devient grise. C'est le
premier défaut de l'étirement décrit dans le rapport : quelques voxels extrêmes, souvent près de bords
nets (graisse très claire à côté de l'air, dont le gradient est amplifié), fixent l'échelle de tout le
volume. Un étirement entre percentiles (par exemple 1\,\% et 99\,\%) serait meilleur ici.

\section{Questions qu'on peut me poser}

\paragraph{Pourquoi $T=60$ et pas 50 ?} 50 est le seuil de l'article pour des photos. Sur l'IRM
normalisée, l'histogramme du centre du cerveau a deux pics : la substance grise vers 120 et la substance
blanche vers 200. Le seuil 60 est nettement en dessous des deux : la zone rehaussée est le liquide
cérébrospinal, l'os cortical et l'air, et le tissu cérébral n'est pas modifié directement. Le choix reste
empirique ; je n'ai pas cherché de seuil optimal.

\paragraph{Est-ce utile cliniquement ?} Pas directement. La méthode ne crée pas d'information : sur le
fantôme, la lésion ressort $2{,}5$ fois plus, mais le bruit aussi, donc le rapport contraste sur bruit ne
change pas. Elle rend des détails sombres plus visibles à l'œil, comme le réglage de fenêtre que font les
radiologues, mais de façon locale (seulement dans la zone sombre) et sans bord visible.

\paragraph{Le solveur tient-il compte des voxels non cubiques ?} Non, et c'est une limite. Le gradient est
calculé par voxel, pas par millimètre : le solveur traite les trois axes de la même façon, alors qu'après
réduction les voxels font $2\times2{,}67\times2{,}67$\,mm. Pour en tenir compte, il faudrait diviser
chaque pulsation par la taille du voxel dans cette direction ($w_m/h_1$, $w_n/h_2$, $w_p/h_3$). Cela
change un peu le poids de chaque direction dans les moindres carrés, mais pas les conclusions.

\paragraph{Pourquoi ne pas traiter le volume complet ?} Pour la mémoire. On pourrait : utiliser une FFT
réelle (\texttt{rfftn}, deux fois moins de mémoire), calculer en simple précision (encore deux fois
moins), ou libérer les tableaux au fur et à mesure. On pourrait aussi remplacer le prolongement miroir et
la FFT par une transformée en cosinus (DCT), qui donne la même condition de Neumann sans doubler la taille.

\paragraph{Pourquoi les stries sont-elles horizontales ?} Parce que j'ai traité les coupes axiales
(perpendiculaires à l'axe vertical $S$). Chaque coupe axiale devient une ligne horizontale sur les vues
sagittale et coronale.

\paragraph{Pourquoi pas de stries sur la coupe axiale ?} Une coupe axiale est une seule coupe traitée en
2D : elle est cohérente avec elle-même. Le défaut est entre les coupes, donc il ne se voit que sur les
vues qui en traversent plusieurs.

\paragraph{D'où vient l'image ?} De OpenNeuro, une base publique de données d'imagerie cérébrale, sous
licence CC0 (domaine public). Elle est anonymisée (visage effacé).

\end{document}
